\section{The instrument: a complete labor account and the workload frontier}
\label{sec:frontier}

\subsection{Count the complete human workload}

Compare a human operating regime $H$ with an agent-assisted regime $A$. For one work package, let
$\lambda$ be visits per period and $h_0>0$ mean baseline case-handling hours. In the agent regime,
$q$ is the fraction requiring substantive human exception handling; $h_E$ is the complete human time
on an exception; $h_R$ is mean human review on a non-exception case; and $h_W$ is additional rework
per visit. Fixed, case-independent human upkeep is $B_H$ or $B_A$.
\begin{align}
\bar h_A &= q h_E + (1-q) h_R + h_W, \label{eq:hbar}\\
L_H &= \lambda_H h_0 + B_H, \qquad L_A = \lambda_A \bar h_A + B_A, \label{eq:hours}\\
N_s &= L_s / K_s, \qquad s\in\{H,A\}. \label{eq:fte}
\end{align}
$N_s$ is workload-equivalent FTE at $K_s$ useful hours per period. Exception effort includes any
preliminary human review on those same cases. $h_R$ includes the expected cost of sampled audits
and mandatory approvals on the ordinary path. $h_W$ excludes work already attributed to these
terms or to repeat visits in $\lambda$. $B_A$ includes policy maintenance, agent evaluation, training,
and incident review within the comparison boundary. Shared teams are allocated once. The
account is invariant to the way an enterprise partitions its gates: the same activities, hours, and
costs give the same total whether they are called one gate or three, so splitting or merging gates
cannot create or hide labor.

This boundary follows the business process rather than a single department's dashboard. A system
that reduces one team's reading time but transfers more remediation to project engineers
has not necessarily reduced enterprise work. Likewise, a human waiting exclusively on an agent
consumes capacity; inference that runs without human attention does not.

\subsection{Separate selection from post-deployment effort}
\label{sec:selection}

A high exception cost has two different explanations: the routed cases were expensive already, or
deployment made them more expensive. Let $h_E^H$ and $h_S^H$ be their conditional baseline means
under a frozen case partition. Define
\begin{equation}
\kappa = \frac{h_E^H}{h_0}, \qquad \sigma = \frac{h_S^H}{h_0}, \qquad
\eta = \frac{h_E}{h_E^H}, \qquad k = \kappa\eta. \label{eq:ratios}
\end{equation}
Here $\kappa$ measures baseline selection; $\eta$ measures post-deployment effort inflation or
assistance for that group. For $0<q<1$ and positive $h_E^H$,
\begin{equation}
q\kappa + (1-q)\sigma = 1, \qquad \sigma = \frac{1-q\kappa}{1-q}\ge 0. \label{eq:partition}
\end{equation}
Consequently $q\kappa\le 1$. This constraint applies to $\kappa$, not to the post-deployment
multiplier $k$. The same baseline exception group can take longer after deployment, so $qk$ may
exceed one. A fixed routing rule can be applied in shadow mode to historical cases to estimate this
partition; changing the rule requires a new partition or explicit reweighting.

\begin{proposition}[Residual selection and retained baseline labor]
\label{prop:selection}
Let $T\ge 0$ denote integrable baseline handling time with $\mathbb{E}[T]=h_0>0$. Let
$J\in\{0,1\}$ indicate assignment to the human residual and $e(T)=\Pr(J=1\mid T)$, with
$q=\mathbb{E}[e(T)]>0$. Then
\begin{align}
\mathbb{E}[T\mid J=1]-h_0 &= \frac{\operatorname{Cov}(T,e(T))}{q}, \label{eq:selection}\\
\phi := \frac{\mathbb{E}[TJ]}{h_0} &= q\kappa = q+\frac{\operatorname{Cov}(T,e(T))}{h_0}. \label{eq:phi}
\end{align}
Thus the residual's baseline labor share exceeds its case share exactly when routing is positively
associated with baseline effort. If $e$ is nondecreasing, this covariance is nonnegative. No ordering
of conditional variances follows.
\end{proposition}

\begin{proof}
The law of iterated expectations gives $\mathbb{E}[TJ]=\mathbb{E}[Te(T)]$. Subtracting $h_0q$ and
dividing by $q$ gives Equation~\eqref{eq:selection}; dividing by $h_0$ instead gives
Equation~\eqref{eq:phi}. For independent copies $T,T'$, the covariance has the representation
\[
2\operatorname{Cov}(T,e(T)) = \mathbb{E}\bigl[(T-T')(e(T)-e(T'))\bigr].
\]
A nondecreasing $e$ makes the integrand nonnegative. The identities involve means and a
cross-moment, not a constraint ordering the variances of $T$ and $T\mid J=1$.
\end{proof}

The economically important quantity is $\phi$, not $q$ alone. A residual containing 20\% of cases can
retain 40\% of baseline labor when $\kappa=2$. With post-deployment inflation $\eta=1.5$, that group's
effort becomes 60\% of the entire baseline case workload before review, rework, and upkeep are
counted. The selection identity supplies the bridge between exception routing and the capacity
calculation. The empirical question in Hypothesis~H3 is whether actual routing produces the positive
covariance.

\subsection{Compression versus expansion}
\label{sec:compression}

For positive baseline means in both groups, define $\mu=h_R/h_S^H$: the fraction of the standard
group's baseline effort required for review. Like $\eta$ for exceptions, $\mu$ measures
post-deployment effort against that group's own baseline. Retain $m=h_R/h_0=\sigma\mu$ for
comparisons in overall-mean units, and set $r=h_W/h_0$, $b_s=B_s/(\lambda h_0)$. At common visits and
useful hours,
\begin{equation}
\rho := \frac{N_A}{N_H}
 = \frac{q\kappa\eta+(1-q)\sigma\mu+r+b_A}{1+b_H}
 = \frac{\phi\eta+(1-\phi)\mu+r+b_A}{1+b_H}. \label{eq:rho}
\end{equation}
The second form reads directly. The share of baseline labor held by the exception group is multiplied
by the change in its effort; the share held by the standard group is multiplied by the fraction of
that effort still spent on review; rework and upkeep are added. $\phi$ can be measured on historical
work before any deployment, whereas $\eta$ and $\mu$ are what the deployment determines.

If $h_S^H=0$, $\mu$ is undefined and the dimensional form~\eqref{eq:hbar} applies. Effort ratios can
exceed one; $q\kappa\le 1$ remains the baseline partition constraint. For fixed $k=\kappa\eta>m$ and
fixed assurance ratios, the equivalent exception-share threshold is
\begin{equation}
q<q^{*}=\frac{1+b_H-m-r-b_A}{k-m}. \label{eq:qstar}
\end{equation}
A negative $q^{*}$ allows no saving at nonnegative $q$; a value above one places the boundary beyond
the feasible case-share range. If $k=m$, the case share has no direct effect. If $k<m$, the inequality
reverses. These are algebraic boundaries, not fitted empirical thresholds.

The same boundary can be read as a margin on exception effort. For $\phi>0$, workload stays below
the baseline only while
\begin{equation}
\eta<\eta^{*}=\frac{1+b_H-(1-\phi)\mu-r-b_A}{\phi}
=\mu+\frac{1+b_H-\mu-r-b_A}{\phi}. \label{eq:etastar}
\end{equation}
Whenever the standard path alone would save labor ($\mu+r+b_A<1+b_H$), the larger the residual's
share of baseline labor, the smaller this margin.

At $q>0$, the boundary is $k^{*}=m+(1+b_H-m-r-b_A)/q$. Figure~\ref{fig:frontier} draws it separately
for each route: fixed upkeep per visit differs with volume, so a single curve cannot classify both.
Workload increases above the applicable curve. A lower escalation rate can thus be outweighed by
costly exception resolution or assurance.

The same equation states what the law of gate substitution requires. Complete substitution drives each weighted
contribution in the numerator of Equation~\eqref{eq:rho} to zero: exception labor, standard-path labor,
rework, and support. The law predicts that endpoint; the equation measures the distance to it, term by
term, and names the term that blocks each route. Tracking the origin of each residual term over time,
which interventions disappeared, which moved, and which are new, is what Hypothesis~H6 tests.

\subsection{Workload, staffing, and cash are distinct outcomes}
\label{sec:staffing}

A workload-equivalent FTE is an effort unit. Staff needed to meet deadlines also depend on arrival
variability, specialization, minimum coverage, and availability of approvers. A utilization constraint
$L\le unK$, with target $u<1$ and staff $n$, is only a starting point; a delay target needs a queueing
or scheduling model. The present calculations hold useful-hour conversion constant, not a measured
waiting-time distribution.

Similarly, saved hours are not automatically avoidable salaries: the component model identifies a
resource change, and a financial counterfactual determines which part becomes cash or additional
output. Section~\ref{sec:workforce} derives the staffing consequence.
